\section{Implementation and Cost Estimation}
\label{sec:implementation}

We describe the system architecture and provide a rigorous cost
estimation for the cryptographic operations. The system is designed as
an extension to the BTV-Transparency stack~\cite{btv-t-repo}.

\subsection{System Architecture}

The implementation consists of three components:
\begin{enumerate}
  \item \textbf{Log Server.} Maintains the Merkle tree state $L_t$.
  Stores Pedersen commitments $C_i$ in the leaves; decision plaintext
  is retained by the operator off-chain and destroyed on redaction.
  \item \textbf{Redaction Service.} Handles authorisation requests,
  generates $\pi_{\mathit{auth}}$ and $\pi_{\mathit{stat}}$, and
  computes the new state $L_{t+1}$. The circuit is implemented in
  Noir~(Aztec) targeting BN254.
  \item \textbf{Verifier/Auditor.} A lightweight client that verifies
  transition proofs against the public log state. Verification is
  constant-time in batch size (Groth16 proof).
\end{enumerate}

\subsection{Circuit Complexity Recap}

As derived in Section~\ref{sec:zk}, the gate count for $N = 1{,}000$
decisions with Merkle depth $d = 20$ and $|\mathcal{G}| = 4$ groups is:

\begin{center}
\begin{tabular}{lrr}
\toprule
Layer & Gates/leaf & Total ($N$=1k) \\
\midrule
Merkle inclusion (Poseidon, $d$=20) & 4{,}400 & 4.4M \\
Pedersen opening (BabyJubJub) & 600 & 0.6M \\
Conditional aggregation (4 groups) & 200 & 0.2M \\
Statistical comparison & $\approx$3 & $<$5k \\
\midrule
Total & & $\approx$5.2M \\
\bottomrule
\end{tabular}
\end{center}

\subsection{Performance Projections}

Based on published benchmarks for the Barretenberg proving backend on
comparable circuits in the 5--10M gate range~\cite{aztec2024}, we
project:

\begin{itemize}
  \item \textbf{Proof generation:} 2.0--2.5~s for $N = 1{,}000$
  (linear in $N$; $\approx$21~s for $N = 10{,}000$).
  \item \textbf{Proof verification:} $<$15~ms, constant in $N$.
  \item \textbf{Proof size:} $\approx$3.2~kB (Groth16; constant in $N$).
\end{itemize}

\paragraph{Transparency statement.}
The figures above are \emph{estimates} derived from the theoretical gate
complexity of the proposed circuit and the known performance
characteristics of the Barretenberg backend for circuits of similar
size. Implementation of the Noir circuit and execution of on-target
benchmarks are ongoing work. We note that even under a pessimistic
$10\times$ overhead (20--25~s per batch), the protocol remains viable
for asynchronous redaction workflows, which are non-interactive and do
not require sub-second latency. Regulators processing GDPR erasure
requests typically operate on timescales of days, not milliseconds.
\subsection{Comparison with Baseline Schemes}

\begin{center}
\begin{tabular}{lccc}
\toprule
Scheme & Selective deletion & Stat.\ integrity & Const.\ verify \\
\midrule
Redactable blockchain~\cite{ateniese2017} & \checkmark & $\times$ & \checkmark \\
Fine-grained rewriting~\cite{derler2019} & \checkmark & $\times$ & \checkmark \\
Accountable Redaction (this work) & \checkmark & \checkmark & \checkmark \\
\bottomrule
\end{tabular}
\end{center}

\noindent The key differentiator is statistical integrity: prior schemes
allow redaction but provide no mechanism to prevent the operator from
exploiting the right to erasure as a selective deletion vector.
